% ============================================================
%  THE ORTYX PROTOCOL
%  Appendix E: Entropic Accessibility and Constraint Descent
% ============================================================

\chapter{Entropic Accessibility and Constraint Descent}
\label{app:entropy}

\noindent
This appendix develops the RSVP-compatible formalism
of accessibility entropy and constraint descent, which
provides the mathematical basis for the reinterpretation
of salience as inverse accessibility and the
constraint-governed trajectory framework of Chapter~17.

\section{Admissible Trajectory Spaces}
\label{app-e:admissible}

Let $X$ be the state space of a dynamical system.
From a state $x_t \in X$ at time $t$, define the
admissible future trajectory space:
\[
  \AdmTraj(x_t) = \{x : X \to X\text{ continuous} :
  x(t) = x_t,\ \mathcal{C}(x(s)) = 0\ \forall s \ge t\},
\]
where $\mathcal{C} : X \to \mathbb{R}_{\ge 0}$ is
a constraint functional with $\mathcal{C}(x) = 0$
denoting constraint satisfaction.

\begin{definition}
The \emph{accessibility entropy} at state $x_t$ is:
\[
  \AccEnt(x_t) = \log\mu\bigl(\AdmTraj(x_t)\bigr),
\]
where $\mu$ is an appropriate measure on the space
of admissible trajectories. High $\AccEnt$ indicates
many admissible futures (low constraint); low $\AccEnt$
indicates few admissible futures (high constraint).
\end{definition}

\section{Functional Systems as Low-Entropy Basins}
\label{app-e:functional}

Functional systems --- systems that reliably achieve
their objectives --- correspond to low-dimensional
admissible basins:
\[
  \mu\bigl(\AdmTraj(x_t)\bigr) \ll
  \mu\bigl(\AdmTraj_{\max}\bigr),
\]
where $\AdmTraj_{\max}$ is the unconstrained trajectory
space. Equivalently, $\AccEnt(x_t)$ is small relative
to its maximum value.

\begin{proposition}
A system in a low-$\AccEnt$ state is highly predictable:
its future trajectory is strongly determined by its
current state. A system in a high-$\AccEnt$ state is
weakly predictable: many distinct futures are consistent
with its current state.
\end{proposition}

The RSVP reinterpretation of salience (\S\ref{app-b:jitter-entropy})
proposes $\SalFunc \propto 1/\AccEnt$: signals with
many admissible continuations are low-salience (high
entropy, unpredictable); signals with few admissible
continuations are high-salience (low entropy, predictable).

\section{Constraint Descent Dynamics}
\label{app-e:constraint-descent}

When the constraint functional $\mathcal{C}$ is
differentiable, constraint-preserving evolution follows:
\[
  \frac{dx}{dt} = -\nabla \AccEnt(x) = -\frac{\nabla\mu(\AdmTraj(x))}{\mu(\AdmTraj(x))},
\]
driving the system toward states with fewer admissible
trajectories --- toward higher constraint and lower
entropy. Stable structures correspond to local minima
of $\AccEnt$:
\[
  \nabla \AccEnt(x^\ast) = 0, \qquad
  \nabla^2 \AccEnt(x^\ast) > 0.
\]

\section{The Salience--Accessibility Correspondence}
\label{app-e:sal-acc}

The structural claim linking the Marine salience metric
to accessibility entropy can be stated as the following
research hypothesis:

\begin{auditprop}
\label{ap:salience-accessibility}
\textit{(Salience--Accessibility Hypothesis)}: For
audio signals, peak-level jitter $\Jp(i) + \Ja(i)$
is negatively correlated with the accessibility entropy
$\AccEnt(x_{t_i})$ of the signal's state at peak $i$,
when the state space $X$ is specified as the space
of admissible waveform continuations under the signal's
physical generating model. This hypothesis generates
a testable prediction: low-jitter peaks should correspond
to states where the waveform's short-term future is
strongly constrained by its current state.
\end{auditprop}

This audit proposition is classified as a \LevelII{}
claim (productive research hypothesis) rather than
a \LevelI{} claim (established engineering result)
because the specification of the state space and measure
required to test it has not been completed.
